% ===================================================================
%  اللوحة رقم ٨ — الحصاد. --- الصيد، قطاف العنب، صيد السمك.
%
%  Traduction, faite en 2026, en arabe levantin d'aujourd'hui, sur
%  l'ido de texto/io. Regles de la colonne : voir 10-lawha-01.tex.
%  Deux scenes, deux numerotations.
%
%  ECART DECLARE, ET IL FAUT LE SAVOIR AVANT D'ECRIRE. Le bloc
%  t08-c2-03-1 est l'un des trois que renvoji.py laisse passer sans
%  rien dire : le fac-simile ido y compose « \textsuperscript{13)} »,
%  malforme, sans parenthese ouvrante, et la machine n'y voit que
%  « 3 12 14 15 16 ». UN ECART DECLARE REND LE BLOC AVEUGLE — le
%  controle n'y regarde plus rien du tout. L'ordre juste, pris sur le
%  volume francais, est 3, 12, 14, 13, 15, 16 ; il a ete pose a la
%  main et recompte a la main.
%
%  CETTE COLONNE EST CHEZ ELLE DANS LA DEUXIEME SCENE, ET DEPUIS PLUS
%  LONGTEMPS QUE PARTOUT AILLEURS DANS LE RELEVE. La vigne a ete
%  domestiquee dans ces montagnes, on y presse du vin depuis six mille
%  ans, et les fouloirs tailles dans le roc sont encore visibles
%  au-dessus des villages de la Bekaa. Rien de ce qui sert au pressoir
%  n'est emprunte : الكرم la vigne (38), المعصرة le pressoir (28),
%  اللولب la vis (29), العنب le raisin (26), القفّة le panier (12),
%  السلافة le premier jus. Et le mot pour le vin, « خمر », est
%  semitique commun, plus vieux que l'arabe lui-meme.
%
%  MAIS LA BIERE DE L'ALINEA 6 EST « البيرة », QUI EST ITALIENNE, ET
%  LE HOUBLON (30) « الجنجل ». SIX MILLE ANS DE VIN, ET LE MOT POUR LA
%  BIERE EST ARRIVE PAR BATEAU. Le releve avait vu la colonne haoussa
%  faire exactement l'inverse au meme alinea : « giya » y designait en
%  propre la biere de mil brassee sur place, et c'est le vin qui avait
%  fallu decrire.
%
%  QUATRE MOTS QUE LA COLONNE GUJARATIE AVAIT DU TRANSCRIRE SONT ICI
%  DES MOTS D'ICI, ET C'EST LA MEME RAISON QU'A LA CHATAIGNE DU
%  TABLEAU 7. Le coquelicot (2) est « شقائق النعمان » — la fleur qui
%  couvre la Galilee et la Bekaa en avril, nommee d'apres un roi
%  lakhmide — la ou le gujarati avait du refuser ખસખસ et ecrire
%  « પોપી ». L'aubepine (22) est « الزعرور », le tilleul (35)
%  « الزيزفون », le chardonneret (36) « الحسّون », l'oiseau en cage de
%  toutes les cours de Damas. Quatre plantes et un oiseau d'Europe
%  tempere qui poussent aussi ici : LE LEVANT EST A CHEVAL SUR LES
%  DEUX FLORES DU LIVRET.
%
%  ET LE CRAPAUD (1) EST « علجوم », DISTINCT DU « ضفدع » DE LA
%  GRENOUILLE (17). La colonne gujaratie n'avait qu'un mot pour les
%  deux et avait du s'en expliquer ; celle-ci les separe comme l'ido
%  les separe. Le meme tableau, deux langues, et la finesse est
%  chaque fois du cote qu'on n'attendait pas.
%
%  Enfin l'arc-en-ciel (52) est « قوس قزح », l'arc de Quzah, qui etait
%  un dieu de l'orage avant l'islam. La colonne haoussa disait « bakan
%  gizo », l'arc de Gizo, le heros fripon des contes. DEUX COLONNES,
%  DEUX ARCS-EN-CIEL, DEUX NOMS PROPRES PLUS VIEUX QUE LA RELIGION
%  D'AUJOURD'HUI.
% ===================================================================

%%K t08-noto-1 noto
\VUnotes{143.2mm}{%
\textit{(*) لأنه هالكلمة مش دقيقة : قطاف كتّان، أو قطاف عنب، أو قطاف
تفّاح ؟ يمكن يكون أحسن نستعمل «} \VUgras{mesono} \textit{» اللي انقترحت
بـ Mondo XIV، صفحة ٢٤٢. بس لهلّق الأكاديمية ما قبلت هالجذر الجديد، وهوّي
عم يستنّى النقاش حسب القاعدة الإيدية المعتادة.}%
}

%%K t08-apar-1 apar
\VUsaut{5.0mm}
\VUtitre{15.0pt}{60}{اللوحة رقم ٨}
\VUsaut{3.0mm}
\VUcentre{13.2pt}{90}{\textit{المشهد الأول.}}
\VUsaut{3.0mm}
\VUpk{10.2pt}{40}{الحصاد (*).}
\VUsaut{4.0mm}
\VUpk{10.2pt}{40}{مناظر الضيعة.}
\VUsaut{3.0mm}

%%K t08-c1-01-1 p
1. --- مع \VUgras{الصيف}، منظر الضيعة تغيّر مرّة تانية. البذار اللي
انزرع بالتلم نبت ؛ ونبتة خضرا زغيرة طلعت من الأرض وعملت سنبلة. والحبّ
تكوّن، وبعدين نضج، ومن هلّق ما ضلّ غير الحصاد.

%%K t08-c1-02-1 p
2. --- \VUgras{زهور الحقل} تفتّحت. \VUgras{زهرة الذرة}
\textsuperscript{(1)}، و\VUgras{شقائق النعمان} \textsuperscript{(2)}،
و\VUgras{قفّاز الثعلب} \textsuperscript{(3)} عم يخلطوا بلون حقول القمح
الواحد تباين \VUgras{تويجاتهن} الطازة المفرح.

%%K t08-c1-03-1 p
3. --- شجر الفواكه لبس ورقه ؛ والفواكه نضجت. و\VUgras{شجرة الكرز}
\textsuperscript{(4)} الزغيرة مليانة \VUgras{كرز} \textsuperscript{(5)} حلو،
والولاد عم يقطفوه وياكلوه بلذّة كبيرة.

%%K t08-c1-04-1 p
4. --- ومن هلّق المواشي فيها تقضّي كل النهار بالهوا الطلق.

%%K t08-c1-04-2 p
\VUgras{الحملان} \textsuperscript{(6)} كبروا وصاروا \VUgras{نعاج}
\textsuperscript{(7)} حلوة و\VUgras{كباش} \textsuperscript{(8)} حلوة، بصوف
كتيف. وعم يرعوا بهدوء \VUgras{عشب} \VUgras{المرعى} \textsuperscript{(9)}
المنقّط بـ\VUgras{الأقحوان}.

%%K t08-c1-04-3 p
وعن قريب رح يبلّشوا يجزّوا \VUgras{هالحيوانات} ليوخدوا
\VUgras{صوفهن} اللي بينعمل منه جوخ تيابنا.

%%K t08-tit-2 sub
\VUsaut{5.0mm}
\VUpk{10.2pt}{40}{الراعي.}
\VUsaut{3.4mm}

%%K t08-c1-05-1 p
5. --- \VUgras{الراعي} \textsuperscript{(10)} مبيّن إنّه مش كتير منتبه
إلهن. بس عم تهمّه العاصفة اللي عم تمرق عند الأفق، و\VUgras{راعي الغنم}
\textsuperscript{(13)}، اللي عم يقرّب \VUgras{مطرته} \textsuperscript{(14)}
لشفايفه، مبيّن أكتر ما بيهتمّ.

%%K t08-c1-06-1 p
6. --- الحرّ خانق ؛ لأنه \VUgras{الكلب} \textsuperscript{(15)} كمان إجا
يقطع عطشه ؛ عم يلحس ميّ \VUgras{الساقية} \textsuperscript{(16)} الباردة وعم
يخوّف \VUgras{ضفدع} \textsuperscript{(17)} زغير مسكين كان مختبّي بعشب
الضفّة. وهاد نطّ بالميّ الصافية وين عم يزهّر \VUgras{زنبق الميّ}
\textsuperscript{(18)}.

%%K t08-c1-07-1 p
7. --- و\VUgras{الذعرة} \textsuperscript{(19)} عم تستحمّي جنب
\VUgras{النبع} \textsuperscript{(20)} اللي عم يفور بين صخرتين، وعم تختبّي
تحت \VUgras{الشوك} \textsuperscript{(21)} و\VUgras{شجيرات الزعرور}
\textsuperscript{(22)}.

%%K t08-tit-3 sub
\VUsaut{5.0mm}
\VUpk{10.2pt}{40}{الحصاد.}
\VUsaut{3.4mm}

%%K t08-c1-08-1 p
8. --- لولاد الضيعة، حصاد القمح وقت كتير مسلّي. عم \VUgras{يشقلبوا}
\textsuperscript{(24)} بفرح على \VUgras{أكوام القش} \textsuperscript{(23)}،
وعم يساعدوا \VUgras{الحصّادين} \textsuperscript{(26)}، ووقت هدول عم يحصدوا
\VUgras{حقل القمح} \textsuperscript{(27)} بـ\VUgras{مناجلهن الكبيرة}
\textsuperscript{(25)} المسنونة منيح على \VUgras{حجر السنّ}
\textsuperscript{(30)}، الولاد عم يجمعوا القمح وعم يربطوه \VUgras{حزم}
\textsuperscript{(31)} أو \VUgras{ربطات} \textsuperscript{(32)}.

%%K t08-c1-09-1 p
9. --- وبعض المرّات بيسمحوا لأكبرهن يقصّ بـ\VUgras{المنجل}
\textsuperscript{(12)} \VUgras{السيقان الدهبية} \textsuperscript{(28)} اللي
شايلة \VUgras{السنابل} \textsuperscript{(29)} التقيلة المليانة حبّ ؛
والزغار، وهنّي عم يحرسوا \VUgras{المواشي} \textsuperscript{(33)} اللي عم
ترعى بـ\VUgras{المرعى} \textsuperscript{(34)} تحت ضلّ \VUgras{الزيزفون}
\textsuperscript{(35)} الكبير، عم يمسكوا بـ\VUgras{شبكتهن}
\textsuperscript{(36)} \VUgras{الفراش} الحلو.

%%K t08-c1-09-2 p
وبعضهن عم يطلعوا على \VUgras{العربايّة} \textsuperscript{(37)} ليرتّبوا
الحزم اللي عم يناولوهن ياها بـ\VUgras{المذراة} \textsuperscript{(38)}.

%%K t08-c1-10-1 p
10. --- وبكل \VUgras{السهل} \textsuperscript{(39)} اللي عم تطير فيه
\VUgras{القبّرات} \textsuperscript{(41)}، في حركة. الكل مشغولين بالحصاد. بس
وقت الفقرا عم يضلّوا يستعملوا العدّة العتيقة — \VUgras{المناجل الكبيرة
والمناجل والنوارج} \textsuperscript{(40)} — أصحاب الملك الأغنيا عم يحصدوا
قمحهن أو \VUgras{جاودارهن} \textsuperscript{(43)} بـ\VUgras{آلة الحصاد}
\textsuperscript{(42)}. وبعدين، بلا ما يحتاجوا يشيلوا الحزم عالبيدر،
بيعملوا بالمحلّ \VUgras{أكداس} \textsuperscript{(44)} كبيرة.

%%K t08-c1-11-1 p
11. --- وبعدين بيجيبوا \VUgras{آلة الدرس} \textsuperscript{(47)} اللي
بتشغّلها \VUgras{المكنة المتنقّلة} \textsuperscript{(45)} بـ\VUgras{سير
النقل} \textsuperscript{(46)}، وهيك الحبّ بينفصل عن \VUgras{التبن} بلا ما
يترك الحقل اللي انزرع فيه.

%%K t08-tit-4 sub
\VUsaut{5.0mm}
\VUpk{10.2pt}{40}{العاصفة.}
\VUsaut{3.4mm}

%%K t08-c1-12-1 p
12. --- كتير مرّات بتصير \VUgras{عواصف} \textsuperscript{(53)} قوية بوقت
الحصاد. أول شي بينسمع هدير \VUgras{الرعد} من بعيد ؛ و\VUgras{ريح الغرب}
\textsuperscript{(54)} القوية بتبلّش تهبّ وبتحني وهيّي عم تزفر أكبر شجر
\VUgras{الحرش} \textsuperscript{(49)}. و\VUgras{الغيوم} \textsuperscript{(56)}
السودا بتهاجم السما ؛ وبتشقّها لمعات \VUgras{البرق} \textsuperscript{(55)}
المتعرّجة اللي بتعمي. و\VUgras{المطر}، وبعض المرّات \VUgras{البرد}،
بيبلّشوا ينزلوا وبيدعسوا الزرع الجاهز للحصاد.

%%K t08-c1-13-1 p
13. --- وكتير مرّات الصاعقة بتحرق بنا. هيك السنة الماضية سقف
\VUgras{طاحون الهوا} \textsuperscript{(50)} صار رماد وتنتين من
\VUgras{جنحانها} \textsuperscript{(51)} تكسّروا.

%%K t08-c1-14-1 p
14. --- وبعد كم ساعة، الهدوء بيرجع. \VUgras{قوس قزح}
\textsuperscript{(52)} لامع بيبيّن عند الأفق والطبيعة بترجع لحياتها اللي
انقطعت كم لحظة. والفلّاحين، اللي التجوا بـ\VUgras{الأكواخ}
\textsuperscript{(48)}، بيرجعوا عالشغل وبيجتهدوا بشجاعة يصلّحوا الضرر
اللي صار.

%%K t08-tit-5 sub
\VUsaut{6.0mm}
\VUcentre{13.2pt}{90}{\textit{المشهد التاني.}}
\VUsaut{3.6mm}

%%K t08-tit-6 sub
\VUpk{10.2pt}{40}{الصيد. --- قطاف العنب.}
\VUsaut{1.4mm}
\VUpk{10.2pt}{40}{صيد السمك.}
\VUsaut{4.0mm}

%%K t08-c2-01-1 p
1. --- بأواخر \VUgras{آب} أو نصّ \VUgras{أيلول}، حسب المناطق، بيفتح
موسم الصيد. وبالكاد أشعّة الشمس الأولى بتطلّع \VUgras{السحالي}
\textsuperscript{(2)} من جحورهن، لحدّ ما \VUgras{الصيّاد}
\textsuperscript{(5)} بيمشي مع \VUgras{كلابه} \textsuperscript{(4)}.

%%K t08-c2-02-1 p
2. --- لبس \VUgras{بدلة الصيد} \textsuperscript{(9)} المتينة من جوخ تخين،
وحطّ برجليه \VUgras{طماقات} جلد، فيها رح يقدر يمشي بالعشب المبلول وين
عم تختبّي \VUgras{العلاجيم} \textsuperscript{(1)} ؛ وأخد \VUgras{بارودته}
\textsuperscript{(6)} و\VUgras{كيس الصيد} \textsuperscript{(10)} تبعه،
وهيّو راح.

%%K t08-c2-03-1 p
3. --- وحماس المطاردة بيوصله لنصّ كرم عم يصير فيه قطاف العنب على
قدم وساق. وولاد صاحب الكرم، اللي بالعادة بيرعوا المعز عالطريق، اليوم
عم يتركوهن يرعوا على كيفهن، وبعد ما ربطوا بوتد \VUgras{تيس}
\textsuperscript{(3)} عتيق ما كان رح يفوّت فرصة حرّيته ليهرب، عم ياخدوا
هنّي كمان نصيبهن من المتعة. واحد منهن عم ياخد عنقود عنب من \VUgras{قفّة
القطاف} \textsuperscript{(12)} وعم يتسلّى وهوّي عم يسيّل \VUgras{العصير}
\textsuperscript{(14)} بـ\VUgras{كباية} \textsuperscript{(13)} ليعمل
\VUgras{سلافة} (نبيد ما اختمر). والتاني عم يحطّ على راسه \VUgras{إكليل
ورق} \textsuperscript{(15)} هلّق ضفره.

وجنبهن، \VUgras{عصافير} \textsuperscript{(16)} قليلة الحيا عم تيجي لحدّ
قدّام المعصرة لتسرق حبّات العنب.

%%K t08-c2-04-1 p
4. --- ووقت الولاد عم يتسلّوا، الرجال عم يشتغلوا. \VUgras{القطّافين}
\textsuperscript{(18)} عم يجيبوا العنب عالقبو بـ\VUgras{القفف تبع الضهر}
\textsuperscript{(17)} ؛ و\VUgras{صانع البراميل} \textsuperscript{(19)} عم
يصلّح \VUgras{البراميل} \textsuperscript{(20)}، وعم يتأكّد إذا
\VUgras{ألواحها} \textsuperscript{(22)} و\VUgras{قاعاتها}
\textsuperscript{(21)} منيحة، وعم يبدّل \VUgras{الأطواق}
\textsuperscript{(23)} المكسورة. وهوّي كمان اللي عمل \VUgras{الأحواض}
\textsuperscript{(25)} الكبيرة اللي بيكبّوا فيها \VUgras{العنب}
\textsuperscript{(26)} ليخمّروه.

%%K t08-c2-05-1 p
5. --- ولمّا بيخلص التخمير، بينقّوا النبيد وبيحطّوا الباقي على
\VUgras{المعصرة} \textsuperscript{(28)}، وبكبس \VUgras{اللولب}
\textsuperscript{(29)} شوي شوي بيعصروا العصير اللي بينزل بـ\VUgras{حوض
زغير} \textsuperscript{(32)}، وبعدها كمان بيطلع \VUgras{نبيد}
\textsuperscript{(31)}.

%%K t08-tit-8 sub
\VUsaut{6.0mm}
\VUpk{10.2pt}{40}{البيرة ونبيد التفّاح.}
\VUsaut{3.4mm}

%%K t08-c2-06-1 p
6. --- وعلى حيط \VUgras{القبو} \textsuperscript{(27)} عم يتسلّق غصن
\VUgras{جنجل} \textsuperscript{(30)}. وهون هوّي بس نبتة زينة، بس ببعض
البلدان، وين الكرمة كتير نادرة، بيزرعوا الجنجل ليعملوا
\VUgras{البيرة}.

%%K t08-c2-07-1 p
7. --- و\VUgras{شجرة التفّاح} \textsuperscript{(33)} الزغيرة محمّلة تفّاح
رح يعمل، لمّا ينضج، \VUgras{نبيد تفّاح} ممتاز. وبـ\VUgras{قفصهن}
\textsuperscript{(34)}، \VUgras{الكناري} \textsuperscript{(35)}
و\VUgras{الحسّون} \textsuperscript{(36)} عم يتطلّعوا بعيون حاسدة عالعصافير
اللي عم تلعب بحرّية.

%%K t08-c2-08-1 p
8. --- ولحدّ ما توصل النظر، على منحدر التلال، مصفوفة \VUgras{دعامات
الكرمة} \textsuperscript{(40)} اللي شايلة \VUgras{جذوع الكرمة}
\textsuperscript{(39)} الحلوة كتير، المحمّلة \VUgras{عناقيد} تقيلة.

%%K t08-c2-08-2 p
وطول السنة، \VUgras{صاحب الكرم} \textsuperscript{(42)} اشتغل ليعتني
بـ\VUgras{كرمه} \textsuperscript{(38)}، وهلّق عم ياخد أجر تعبه بقطاف حلو.
و\VUgras{القطّافات} \textsuperscript{(41)} عم يعبّوا الأحواض المحطوطة
عالعربايّة اللي عم تجرّها \VUgras{بغال} \textsuperscript{(43)}، والكل
مبسوطين.

%%K t08-tit-9 sub
\VUsaut{5.0mm}
\VUpk{10.2pt}{40}{صيد السمك.}
\VUsaut{3.4mm}

%%K t08-c2-09-1 p
9. --- ونفس الشي عند \VUgras{صيّادي السمك بالسنّارة} \textsuperscript{(44)}
اللي من الصبح إجوا قعدوا على ضفّة الميّ مع \VUgras{قصبات صيدهن}
\textsuperscript{(45)}.

%%K t08-c2-09-2 p
و\VUgras{السمك} \textsuperscript{(47)} بيعضّ \VUgras{الصنّارة} أول ما
يرموا \VUgras{خيطهن} \textsuperscript{(46)} بالميّ، وبالكاد بيلحقوا
يجدّدوا \VUgras{الطعم} لحدّ ما ضحية جديدة بتترفرف بطرف الخيط.

%%K t08-c2-09-3 p
ومع هيك \VUgras{الصيّاد بالشبكة} \textsuperscript{(48)}، اللي من
\VUgras{مركبه} \textsuperscript{(49)} عم يرمي \VUgras{الشبكة} بنصّ
\VUgras{النهر} \textsuperscript{(50)}، عم يمسك سمك أكتر بكتير.

%%K t08-c2-10-1 p
10. --- \VUgras{الخريف} موسم اللفّات الحلوة : مشي، \VUgras{عالخيل}
\textsuperscript{(52)}، \VUgras{عالبسكليت} \textsuperscript{(53)}،
\VUgras{عالسيّارة} \textsuperscript{(54)}. و\VUgras{الطرقات}
\textsuperscript{(51)} نضيفة والناس عم تستفيد من آخر الأيام الحلوة، لأنه
هيّو من هلّق \VUgras{السنونو} \textsuperscript{(56)} عم يتجمّع عالسطوح
ليطير بكل جناحيه لبلدان أدفا. وورق \VUgras{شجرة الجوز}
\textsuperscript{(57)} العتيقة عم يصفرّ، و\VUgras{الجوز} عم يوقع،
و\VUgras{الشجر} العالي المتجمّع متل \VUgras{باقة} \textsuperscript{(58)}
عالـ\VUgras{مرتفع} \textsuperscript{(59)} عن قريب رح يفقد ورقه.

%%K t08-c2-11-1 p
11. --- \VUgras{الشمس} \textsuperscript{(60)} عم تطلع أتأخّر، والنهار عم
يقصر، وبرد لاذع شوي عم يبلّش ينحسّ عالـ\VUgras{صبح}، وهيّي من هلّق
أسراب \VUgras{الشنقب} \textsuperscript{(61)} عم تترك بلادنا اللي ما عاد
فيها ضيافة.

\VUsaut{6.0mm}
\VUfilet{22mm}
